% Ad Atticum 2.18 --- headnote
% ad-atticum-02-18, June or Quintilis (July) 59 BC, written at Rome

\textit{Cicero to Atticus, June or July 59 BC. Written from
Rome in the high triumviral summer, when the Julian
agrarian laws have passed and the political mood has set
hard. \textsection 1 names the one man still openly
resisting in the Forum: the young Curio (the elder son of
the consul of 76, the future tribune of 50 BC and Caesarian
partisan), in this season the popular optimate hope. The
shouts that follow him through the Forum are matched by the
hisses that pursue Fufius the Caesarian tribune. \textsection
2 turns to the Campanian law: the new clause swearing
candidates to the agrarian settlement has just been passed,
and most candidates take the oath; only Laterensis, withdrawing
from the tribunate, refuses. \textsection 3 is the personal
crux of the letter: Caesar has offered Cicero a legateship
under him in Gaul, and a \textit{libera legatio} is also on
offer; Cicero is keeping his options open but does not yet
intend to take either, since he wishes to fight rather than
flee. The threat is named: \textit{Pulchellus}, ``little
Pulcher'' --- Clodius, with the diminutive that recurs in
this stretch of letters. \textsection 4 closes with the
manumission of Statius, Quintus's freedman, which has cost
Cicero some grief, and the brother's plea: hold yourself
ready to fly when I call.}
